\section{失败、长视频与 Q\&A 的工程边界}

能力和 scaling 之后，最后必须回到失败。复杂碰撞揭示模型缺乏稳定多体交互；长视频受到 token 与二次 attention 制约；reasoning 需要媒体形式的 trace 与 verifier；部署还面对 watermark、数据污染和未公开 serving architecture。本章把 Q\&A 中的口头边界融入主线。

\subsection{复杂交互为什么仍会崩}

评测与 scaling 说明模型在平均意义上很强，本节故意转向最不平均的部分：复杂交互失败。碰撞样例不是偶然笑料，而是一个诊断；当多个对象高速接触、遮挡并改变拓扑关系时，局部视觉 plausibility 不足以保证全局状态一致。读图时看车体是否穿透、碎片来源是否连续、碰撞前后身份是否保持。

\lecturefigure{slide-460-complex-motion-failure.jpg}{失败案例：复杂车辆碰撞中出现几何和交互不一致}{00:50:04}

\begin{warningbox}{视频失败常是状态管理失败}
模型可能逐帧生成“看起来像碰撞”的纹理，却没有稳定对象状态、接触约束和事件顺序。提高单帧清晰度无法自动修复这类错误；需要更好的时空表示、数据、长期记忆、约束或自校正机制。
\end{warningbox}

\subsection{长视频首先是序列长度问题}

设每秒产生 $n_f$ 个 latent tokens，时长为 $S$，则 $N\approx n_fS$。若使用全注意力，主项随时长近似
\[
\operatorname{cost}(S)=O((n_fS)^2d).
\]
把 16 秒扩为 32 秒，在其他条件不变时 token 数约翻倍，attention 主项约四倍；还要保持更长情节、角色和相机状态。

\begin{knowledgebox}{长视频的三层困难}
第一层是计算：context、激活和通信增长；第二层是数据：长视频 caption 与事件边界更难获得；第三层是建模：模型要维持身份、因果链和叙事目标。仅靠更大显存只能缓解第一层。
\end{knowledgebox}

\teachervoice{00:57:11--01:00:13，老师说 73K 已接近当时可行训练长度。若时长翻倍，序列与 attention 成本迅速增加；长电影还需要超越朴素一次性全序列生成的系统设计。}

\subsection{Reasoning、verifier 与 native media generation}

前一节已经定位复杂交互的状态管理错误，本节追问系统能否主动发现并修正它。讲者提出两个开放方向：其一，让视频模型在输出前发现明显错误并自校正；其二，把视频生成纳入原生多模态 LLM。两者都需要回答学习目标是否统一，以及如何验证媒体输出“正确”，而这比增加模型参数更难定义。

\lecturefigure{slide-464-future-directions.jpg}{未来方向：继续 scaling、媒体 reasoning 与原生多模态生成}{00:53:37}

若把生成--检查--修正写成循环，可抽象为
\[
y^{(k)}=G_\theta(z,c,r^{(k)}),\qquad
e^{(k)}=V_\phi(y^{(k)},c),\qquad
r^{(k+1)}=R(r^{(k)},e^{(k)}),
\]
其中 $G$ 是生成器，$V$ 是 verifier，$R$ 更新 reasoning state。困难在于视频没有像数学题那样便宜、确定的正确性判据；人类能一眼发现撞车不合理，并不等于存在可训练的自动 reward。

\begin{importantbox}{媒体 reasoning 的核心瓶颈是验证}
生成“思考文本”并不自动改善视频。必须定义可检测的错误：对象身份、动作顺序、几何、物理、文本约束或审美，并证明 verifier 与人类判断相关。否则 reasoning trace 可能只是额外 token，而非有效自校正。
\end{importantbox}

\subsection{Q\&A 的五条校正}

最后的问答把架构、数据与治理边界说得更清楚。下面五条不是附录，而是理解主讲结论的必要条件。

\begin{enumerate}
  \item \textbf{Inductive bias：}CNN 在小规模或特定机器人 diffusion policy 上可能更有优势；Transformer 更容易沿成熟基础设施扩展，但辩论没有结束。
  \item \textbf{Physics：}可硬编码先验、加入模拟数据或依赖 scale，哪条路线最好仍是开放问题；当前样例不能证明物理已解决。
  \item \textbf{Synthetic data：}低质量生成内容进入 pretraining 会污染数据，但合成数据并非天然有害；受控 post-training 本就常用模型生成样本。
  \item \textbf{Safety：}deepfake 与 watermarking 是独立研发和部署问题，不能由生成质量指标替代。
  \item \textbf{Public evidence：}论文公开训练 infra 细节，课堂没有公开完整 serving architecture；不应猜测单次请求使用多少 GPU。
\end{enumerate}

\begin{knowledgebox}{“统一”应放在哪一层}
Q\&A 中老师指出，只要新模态能编码为 token，Transformer core 可以复用；真正困难往往在表示和目标。统一主干并不要求统一 tokenizer、decoder、loss 或评估协议。
\end{knowledgebox}

\teachervoice{00:54:10--00:56:03，老师承认专用 CNN 的归纳偏置在小规模可有价值；Movie Gen 的经验是大规模 Transformer 更容易扩展、基础设施更成熟，但他明确说自己没有为争论盖棺定论。}

\teachervoice{01:09:23--01:10:33，关于合成数据污染，老师区分 pretraining 中的低质量假视频与受控 post-training 数据：生成数据不是天然坏，关键是质量、目标和过滤。}

\subsection{本章小结}

复杂交互失败揭示视频模型的状态管理不足；长视频同时受二次计算、数据与长期一致性限制；reasoning 的关键是 verifier，而非额外文字。Q\&A 进一步校正了 Transformer、物理先验、合成数据、安全和公开系统证据的边界。

